% Custom title block for Paper 03: Semantic Orbital Mechanics

\begingroup
\setlength{\parindent}{0pt}
\begin{center}
{\LARGE\bfseries Semantic Orbital Mechanics:\\Measuring and Guiding AI Conversation Dynamics\par}
\vspace{1.5em}
{\large Juan Jacobo Jimenez Sanchez\par}
\vspace{0.3em}
{\normalsize Aicoevolution Ltd\par}
{\normalsize Auckland, New Zealand\par}
{\normalsize \texttt{research@aicoevolution.com}\par}
\vspace{1em}
{\normalsize January 2026\par}
\end{center}
\endgroup

\vspace{1.5em}

\noindent\textbf{Abstract}

\vspace{0.5em}

\noindent Current alignment techniques focus largely on response optimization---ensuring individual model outputs match human preferences. However, alignment is fundamentally a dynamical process: meaning emerges not from isolated tokens but from the trajectory of interaction over time. In this paper, we introduce a physics-inspired framework for measuring and guiding these dynamics, treating conversation as an orbital system in high-dimensional semantic space.

\vspace{0.5em}

\noindent Building on the geometric verification of the S64 symbolic framework (Paper 02), we adopt the Semantic Grounding Index (SGI) from Mar\'in's geometric hallucination detection work and reinterpret it as an orbital radius that measures the tension between local responsiveness (query gravity) and global context (history gravity). We define the Conversational Coherence Region as a stable orbit where exploration and grounding are balanced. We then introduce the Semantic Transducer---a telemetry system that decomposes embedding trajectories into actionable S64 signals: symbols, paths, and transformation phases.

\vspace{0.5em}

\noindent To validate this framework, we conduct a controlled steering experiment. By injecting fake telemetry metrics into an AI's system prompt, we test whether conversational orbits can be predictably altered. The results are surprising: conversations maintain orbital stability despite one participant's distorted perception. The orbit is robust; steering affects what is discussed but not how meaning moves.

\vspace{0.5em}

\noindent This finding reveals both the power and limits of orbital dynamics. The transducer provides reliable telemetry---trajectory metrics are invariant across 10 embedding backends. But orbital mechanics describes only the \emph{horizontal} plane of conversation: position and velocity. Detecting semantic manipulation requires the \emph{vertical} dimension---symbolic depth, transformation richness, contribution asymmetry---that reveals not just where meaning is but how deep it goes. This paper reframes the 3-Body Problem of human-AI interaction as a measurable dynamical system, while acknowledging that complete navigation requires the vertical instruments of Paper 04.

\vspace{1em}

\noindent\textbf{Keywords:} semantic orbital mechanics, conversation dynamics, orbital robustness, AI alignment, S64 framework, high-dimensional geometry, semantic transducer

\clearpage

